\cotefile{161-1}
\scanpage{4}

% loc: upper right
\margin{Foncteurs adjoints\\[2pt] [demander Schanuel]}

% loc: upper left
Soit

% loc: upper centre
\diagram{\begin{tikzcd}[ampersand replacement=\&,column sep=large]
E\arrow[r,shift left=.7ex,"u"]\& F\arrow[l,shift left=.7ex,"v"]\&[-.5em]
\begin{aligned}i&:\operatorname{id}_{E}\longrightarrow vu\\ j&:uv\longrightarrow\operatorname{id}_{F}\end{aligned}
\end{tikzcd}}

% loc: below diagram
un couple de foncteurs adjoints.

% loc: a, opening
a) Soit $E'$ une sous-catégorie \deleted{\illegible{1}} pleine de $E$, \deleted{\unsure{et stable par}} Notons que pour $X,Y\in\operatorname{Ob}E'$, l'application

% loc: a, Hom display
\[
\begin{aligned}
\operatorname{Hom}_{E'}(X,Y)&\longrightarrow\operatorname{Hom}_{F}(uX,uY)\quad(\simeq\operatorname{Hom}(X,vuY))\\[-2pt]
\rotatebox{90}{$\simeq$}&\\[-2pt]
\operatorname{Hom}_{E}(X,Y)&
\end{aligned}
\]

% loc: a, explanation
est déduite de l'application $i(Y):\deleted{\unsure{vu}}Y\longrightarrow vu(Y)$ comme $\operatorname{Hom}(X,i(Y))$. Donc \deleted{\unsure{dire}}\inserted{pour} que \deleted{\unsure{ceci}}\inserted{ce soit injectif} \deleted{\unsure{signifie}}\inserted{(resp.)} pour $\forall X,Y\in\operatorname{Ob}E'$ \deleted{\unsure{signifie}}\inserted{i.e. $u|E'$ \deleted{pl.} fidèle (resp. pl. fidèle)} il suffit que $\forall Y\in\operatorname{Ob}E'$ \deleted{\unsure{on ait}}, $i(Y):\deleted{\unsure{soit}}Y\to vu(Y)$ soit \deleted{\unsure{un isom.}}\inserted{un mono} (resp. un \uline{\unsure{isom|mono}}). C'est aussi nécessaire si on sait que $vu(E')\subset\overline{E'}$ (image ess. de l'inclusion $E'\subset\unsure{\overline E|E}$\unsure{)}.

% loc: a, transition
\deleted{Donc Donc \unsure{l'on a}} On trouve donc

% loc: a, condition 1
\[
u|E'\ \text{pl. fidèle et }vu(E')\subset\overline{E'}\quad\Longleftrightarrow\quad
\forall\,Y\in\operatorname{Ob}E',\quad i(Y):Y\xrightarrow{\sim}vu(Y)\tag*{\textcircled{1}}\]

% loc: under condition 1
\margin{donner un exemple montrant que cette condition n'est pas surabondante, avec $E'$ réduit à un élément}

% loc: b, opening
b) Soit $F'$ une sous-catégorie pleine de $F$. Notons que pour $P,Q\in\operatorname{Ob}F'$, l'app.

% loc: b, Hom display
\[
\begin{aligned}
\operatorname{Hom}_{F'}(P,Q)&\longrightarrow\operatorname{Hom}(vP,vQ)\quad(\simeq\operatorname{Hom}(uvP,Q))\\[-2pt]
\rotatebox{90}{$\simeq$}&\\[-2pt]\operatorname{Hom}_{F}(P,Q)&
\end{aligned}
\]

% loc: b, explanation
est déduite de $j(P):uv(P)\longrightarrow P$ par $\operatorname{Hom}(j(P),Q)$. Donc on trouve pour \inserted{pour que} $v|F'$ soit fidèle \inserted{(pl. fidèle)}, il suffit que $\forall P\in\operatorname{Ob}F'$, $j(P):\inserted{uv(P)\to P}$ soit un épi. (resp. un isom). Donc on trouve

% loc: b, condition 2
\[
v|F'\ \text{pl. fidèle et }uv(F')\subset\overline{F'}\quad\Longleftrightarrow\quad
\forall P\in\operatorname{Ob}F',\ j(P):uv(P)\xrightarrow{\sim}P.\tag*{\textcircled{2}}\]

% loc: below condition 2, barred two-line draft
\deleted{Si $\overline{F'}$ est stable par quotients, on a}\\
\deleted{$v|F'$ fidèle et $uv(F')\subset\overline{\unsure{F|F'}}\Longleftrightarrow\forall P\in\operatorname{Ob}F',\ j(P):uv(P)\to P$ un}

% loc: c, correspondence
c) L'ensemble des $E'$ \inserted{strictement pleines $\subset E$} satisfaisant (1), et des $F'$ \inserted{strictement pleines} satisfaisant (2), sont en correspondance biunivoque, en posant

% loc: c, images
\[
\begin{aligned}
\alpha(E')&=\text{image essentielle de }u|E'=\overline{u(E')}\\
\beta(F')&=\text{image essentielle de }v|F'=\overline{v(F')}.
\end{aligned}
\]

% loc: c, equivalences
Si $E'$ et $F'$ \deleted{\unsure{se correspondent ainsi}, alors \unsure{$u|E'$ et $v'$}}\inserted{sont des sous-catégories \inserted{str.} pleines \deleted{\unsure{l'une l'autre}} de $E$ resp. $F$}, les conditions suivantes sont équivalentes :

% loc: c, lines 1 to 5
1) $E'$ satisfait \textcircled{1}, $F'$ satisfait \textcircled{2}, et $E'$ et $F'$ se correspondent comme dessus

2) $u(E')\subset F'$, $v(F')\subset E'$, et $u'$ \deleted{est pl. fidèle}\inserted{et $v'$} sont pl. fidèles

3) \rule{1.1em}{.4pt}\quad\rule{2.1em}{.4pt}\quad et $u'$ une équivalence

4) \rule{3.8em}{.4pt}\quad et $v'$ une équivalence

5) \rule{3.8em}{.4pt}\quad et $u'$ et $v'$ pl. fid.

% loc: c, right of last lines
$\unsure{6|5})\ \ldots\ldots$, et $\left\{\begin{aligned}
\forall\,\deleted{\unsure{X}}Y\in\operatorname{Ob}E',\ i(Y)&:Y\xrightarrow{\sim}vuY\\
\forall\,P\in\operatorname{Ob}F',\ j(P)&:uvP\xrightarrow{\sim}P
\end{aligned}\right.$

% archive sheet number: 3
\scanpage{5}

% loc: d, upper part
d) \deleted{\unsure{Prenons donc maintenant|Passons maintenant}} Soit $E_0$ \deleted{\unsure{=}} la sous catégorie \inserted{str.} pleine de $E$ formée des $Y\in\operatorname{Ob}E$ tels que $i(Y):Y\xrightarrow{\sim}vu(Y)$. C'est donc la plus grande des sous-catégories pleines $E'$ de $E$ satisfaisant (1), et si $E'$ est une ss-catégorie pleine de $E$, elle satisfait (1) sss $E'\subset E_0$. De m\textsuperscript{e} on définit la sous-catégorie \inserted{str.} pleine $F_0$ de $F$ \inserted{formée} des \deleted{$\unsure{Q}$} $P\in\operatorname{Ob}F$ tels que $j(P):uv(P)\xrightarrow{\sim}P$, une sous catégorie pleine $F'$ de $F$ satisfait (2) sss $F'\subset F_0$. Ceci posé, on a $F_0=\overline{u(E_0)}$, $E_0=\overline{v(F_0)}$ (i.e. $F_0$\inserted{$=$} image ess. de $u|E_0$, $E_0=$ image ess. de $v|F_0$) et $u$ et $v$ induisent sur $E_0$ et $F_0$ des équivalences inverses l'une de l'autre.

% loc: e, opening
e) Pour que \deleted{l'image essentielle} l'on \deleted{ait} $v(F)\subset E_0$, il f. et s. que pour $\forall$ \deleted{$\unsure{Q}$} $P\in\operatorname{Ob}F$, on ait \deleted{$\unsure{v(P)\to vuv(P)}$}

% loc: e, isomorphism and return arrow
\diagram{\begin{tikzcd}[ampersand replacement=\&,column sep=huge]
v(P)\arrow[r,"{i(v(P))}","{\sim}" description]\&vuv(P)\arrow[l,bend left=60,dashed,"{v(j(P))\ \text{sera alors l'inverse ??}}"']
\end{tikzcd}}

% loc: e, right margin
\margin{Ou encore, il f. et s. que $v$ se factorise en $F\xrightarrow{v_1}E_1\xrightarrow{\alpha}E$, avec $\alpha$ pl. fidèle, $v_1$ \unsure{passage|pass.} à cat. des fractions, [\uline{NB} revient à $v_1$ avec un adjoint à g. $u_1=u|E_1$ mais pas néc. $\alpha$]}

% loc: e, first left diagram
\diagram{\begin{tikzcd}[ampersand replacement=\&] E\arrow[r,"u"]\arrow[d]\&F\arrow[d]\arrow[dl,"v'"']\\E_0\arrow[r,shift left=1ex,"u_0" description,"\simeq"]\&F_0\arrow[l,shift left=1ex,"v_0"'," \simeq" description]\end{tikzcd}}

% loc: e, consequence
Dans ce cas, $E_0$ est l'image \deleted{\unsure{ess.}}\inserted{essentielle} de $v$, et l'inclusion $F_0\subset\underset{\alpha}{F}$ admet un adjoint à \deleted{\unsure{gauche}}\inserted{droite}, qui est \deleted{$\unsure{une restriction à}\ F_0$} $u_0\circ\underbrace{(v\ \text{astreint à }\unsure{E|E_0})}_{v'}$ :

% loc: e, cancelled opening of calculation
\deleted{Pour que l'on \unsure{ait}}\\
$\deleted{\operatorname{Hom}_{E_0}(X_0,v'Y)\simeq\operatorname{Hom}_{E}(X_0,vu(Y))\simeq\operatorname{Hom}_{F}(uX_0,u(Y))}$

% loc: e, cancelled diagram in left margin
\diagram{\begin{tikzcd}[ampersand replacement=\&]\deleted{E}\arrow[r,"\deleted{u}"]\arrow[d]\arrow[dr,"\deleted{\unsure{u'}}" description]\&\deleted{F}\arrow[d]\arrow[dl]\\\deleted{E_0}\&\deleted{F_0}\end{tikzcd}}

% loc: e, crossed block of Hom identities
\[
\begin{aligned}
&\deleted{\operatorname{Hom}_{F_0}(u_0v'(\underset{\in\operatorname{Ob}F}{P}),\underset{\in\operatorname{Ob}F_0}{Q_0})=\operatorname{Hom}_{F}(uv(P),Q_0)
 \underset{\text{adjonction}}{\simeq}\operatorname{Hom}_{E}(vP,vQ_0)}\\
&\deleted{\simeq\operatorname{Hom}_{E}(\underbrace{vuv(P)}_{vu_0v'(P)},v(Q_0))}\\
&\deleted{\text{car }vP\simeq vuvP\longleftarrow\ \unsure{\text{redonne}}}\\
&\deleted{\simeq\operatorname{Hom}_{E_0}(v_0u_0v'P,v_0(Q_0))}\\
&\deleted{\simeq\operatorname{Hom}_{F_0}(u_0}
\end{aligned}
\]

% loc: e, proof in brackets
[En effet, on a pour $Q_0\in\operatorname{Ob}F_0$, $P\in\operatorname{Ob}F$

% loc: e, proof display
\[
\begin{aligned}
\operatorname{Hom}_{F}(Q_0,P)&\underset{\text{car }u_0v_0\simeq\operatorname{id}}{\simeq}
\operatorname{Hom}_{F}(u_0v_0Q_0,P)\simeq\operatorname{Hom}_{E}(v_0Q_0,vP)=\operatorname{Hom}_{E_0}(v_0(Q_0),v'P)\\
&\underset{u_0\ \text{pl. fid}}{\simeq}\operatorname{Hom}_{F_0}(u_0v_0Q_0,u_0v'P)
\underset{(u_0v_0\simeq\operatorname{id})}{\simeq}\operatorname{Hom}_{F_0}(Q_0,u_0v'P)\big].
\end{aligned}
\]

% loc: lower left, reverse diagram
\diagram{\begin{tikzcd}[ampersand replacement=\&]E\arrow[d]\arrow[dr,"u'"]\&F\arrow[l,"v"']\arrow[d]\\ E_0\arrow[r,shift left=.8ex,"u_0","{\simeq}" description]\&F_0\arrow[l,shift left=.8ex,"v_0"'," {\simeq}" description]\end{tikzcd}}

% loc: lower bracket
[\unsure{De m\textsuperscript{e}}, \deleted{on trouve}\inserted{est-il vrai} que $E_0\hookrightarrow E$ admet $v'u$ comme adjoint à \deleted{\unsure{gauche}}\inserted{gauche}, ou qu'il admet un adjoint à \deleted{\unsure{gauche}}\unsure{?} ?? Revient à ceci : si $\alpha v_1$ a un adjoint à g., avec $\alpha$ pl. fid et $v_1$ passage à cat. fractions, a-t-on que $\alpha$ a un adj. à g ? ]

% loc: bottom, cancelled calculation
\[
\deleted{\operatorname{Hom}_{E}(\underset{\operatorname{Ob}E}{Y},\underset{\operatorname{Ob}E_0}{X_0})\underset{v_0u_0\simeq\operatorname{id}}{\simeq}\operatorname{Hom}_{E}(Y,v_0u_0X_0)\underset{\text{adjonction}}{\simeq}\operatorname{Hom}_{\unsure{E|F}}(uY,u_0X_0)}
\]

% loc: foot, continuation
De façon duale, pour que l'on ait $u(E)\subset F_0$, il f. et s. que l'on ait \deleted{pour} $\forall X\in\operatorname{Ob}E\quad uvu(X)\xrightarrow[\sim]{j(u(X))}u(X)$, et alors $F_0$ est l'image ess. de $E$ par

% archive sheet number: 4
\scanpage{6}

% loc: top, continuation of e
$u$. Dans ce cas, l'inclusion $E_0\hookrightarrow E$ admet un adjoint à \uline{gauche}, savoir $v_0\circ\underbrace{(u\ \text{astreint à }F_0)}_{u'}$ :

% loc: top, Hom calculation
\[
\begin{aligned}
\operatorname{Hom}_{E}(\underset{\in\operatorname{Ob}E}{Y},\underset{\in\operatorname{Ob}E_0}{X_0})
&\underset{v_0u_0\simeq\operatorname{id}}{\simeq}\operatorname{Hom}_{E}(Y,v_0u_0X_0)=\operatorname{Hom}_{E}(Y,vu_0X_0)\underset{\text{adjonction}}{\simeq}\operatorname{Hom}_{F}(uY,u_0X_0)\\
&=\operatorname{Hom}_{F_0}(u'Y,u_0X_0)\underset{v_0\,\text{pl. fid}}{\simeq}\operatorname{Hom}_{E_0}(v_0u'Y,v_0u_0X_0)\\
&\underset{v_0u_0\simeq\operatorname{id}}{=}\operatorname{Hom}_{E_0}(v_0u'Y,X_0)\qquad\text{cqfd.}
\end{aligned}
\]

% loc: f, opening
f) Conditions équivalentes :

% loc: f, items 1 and 2
1) $\operatorname{Im}v\subset E_0$ et \inserted{$1')$} $\operatorname{Im}u\subsetE_0$

2) $v\xrightarrow{\sim}vuv$ et \inserted{$2')$} $u\xleftarrow{\sim}uvu$

% loc: f, bracketed items 3 and 3 prime
[3) \deleted{$u$ se factorise $u$ se factorise en un composé} $\operatorname{Im}v\subset E_0$ \uline{et} l'inclusion $E_0\hookrightarrow E$ admet un adjoint à g

$3')$ $\operatorname{Im}u\subset F_0$ et l'inclusion $F_0\hookrightarrow F$ admet un adjoint à dr. ]

% loc: f, beside items 4 and 4 prime
% slanted
\margin{à voir si on peut simplifier}

% loc: f, item 4
4) $u$ se factorise en $E\xrightarrow{u_1}F_1\overset{\beta}{\unsure{\longrightarrow|\longleftarrow}}F$, avec $\beta$ pleinement fidèle, $u_1$ \deleted{\unsure{un épigraphe}, et $\beta$ et $u_1$ ont des adjoints à droite}\inserted{passage à catégorie des fractions [, $\beta$ et $u_1$ ayant des adjoints à dr. ]}

% loc: f, item 4 prime
$4')$ $v$ se factorise en $F\overset{v_1}{\unsure{\longrightarrow|\text{---}}} E_1\overset{\alpha}{\unsure{\longrightarrow|\text{---}}}E$, avec $\alpha$ pl. fidèle, $v_1$ passage à cat. des fractions, [ $v_1$ et $\alpha$ admettent un adj. à g. ]

% loc: lower paragraph
De \deleted{\unsure{5)|4)}} plus, la donnée d'un couple de foncteurs adjoints $(u,v)$ entre deux catégories $E,F$, satisfaisant les conditions équivalentes 1) à 4), équivaut à la donnée d'une sous-cat. str. pleine $E_0$ de $E$ telle que $E_0\overset{\alpha}{\hookrightarrow}E$ admette un adjoint à g \inserted{$\alpha'$}, d'une sous-cat. str. pleine $F_0$ de $F$ telle que $\beta:F_0\hookrightarrow F$ admette un adjoint à droite \inserted{$\beta'$}, et d'une équivalence entre $E_0$ et $F_0$

% loc: lower diagram
\diagram{\begin{tikzcd}[ampersand replacement=\&] E_0\arrow[r,shift left=.7ex,"u_0"]\&F_0\arrow[l,shift left=.7ex,"v_0"]\end{tikzcd}}

% loc: lower closing
en posant $u=\deleted{\unsure{\beta}}\beta u_0\alpha'$, $v=\alpha v_0\beta'$, les données d'adjonction étant évidentes.

% archive sheet number: 5
\scanpage{7}

% loc: top left
\diagram{\begin{tikzcd}[ampersand replacement=\&]E\arrow[r,shift left=.7ex,"u"]\& F\arrow[l,shift left=.7ex,"v"]\end{tikzcd}}

% loc: upper example a
a) $E$ catégorie ponctuelle, $v$ l'unique foncteur $F\to E$. Il a un adjoint à g. \inserted{$u$} sss $F$ a un objet initial, et alors $u(P)=\varnothing_F$. Alors $u$ est pl. fidèle.

% loc: upper example b
b) $F$ catégorie ponctuelle, $u$ l'unique foncteur $E\to F$. Alors l'adjoint à droite $v$ existe sss $E$ a un objet final, $v(P)=e_E$. Alors $v$ est pl. fidèle.

% loc: left of main paragraph
\diagram{\begin{tikzcd}[ampersand replacement=\&,row sep=normal]E\arrow[d,bend right=50,"{(\alpha')}"']\\E_1\arrow[u,hook',"{\alpha\ \text{pl. fid}}"']\end{tikzcd}}

% loc: middle opening
On suppose que $\alpha$ admet un adjoint à \uline{gauche} $\alpha'$, d'où couple $\alpha',\alpha$ ; \deleted{qui soit} de sorte que $\alpha'$ est un passage à une catégorie de fractions. Considérons alors un autre foncteur pl. fidèle $E_1\xrightarrow{\beta} F$, et supposons que $u=\beta\alpha'$ :

% loc: small triangle
\diagram{\begin{tikzcd}[ampersand replacement=\&] E\arrow[rr,"u"]\arrow[dr,"\alpha'"']\&\&F\\\&E_1\arrow[ur,"\beta"']\&\end{tikzcd}}

% loc: middle question
ait un adjoint à \deleted{gauche}\inserted{droite} (cela exige de toutes façons, sans savoir d'où provient $\alpha'$, que $\alpha$ ait un adjoint à droite \inserted{$\unsure{\alpha|\gamma}$}, \deleted{\unsure{et ce}} ici le fait que $\alpha$ pl. fidèle équivaut à $\alpha'$ passage à cat. des fractions !) Est-il vrai alors que $\beta$ ait un adjoint à droite ? On peut supposer que \deleted{$\unsure{\operatorname{Ob} F}$} $\operatorname{Ob}F=\operatorname{Ob}E_1\cup\{a\}$, il est donné alors par deux foncteurs $\varphi:E_1\to(\mathrm{Ens})$, $\psi:\deleted{\unsure{(\mathrm{Ens})}}E_1^{\circ}\unsure{\to|{-}}(\mathrm{Ens})$, savoir

% loc: pair of functors
\[\varphi(X)=\operatorname{Hom}(a,\beta X)\qquad\psi(X)=\operatorname{Hom}(\beta X,a)\]

% loc: before product
et un morphisme bifonctoriel

% loc: product
\[\unsure{\varpi|\omega}\quad\deleted{\operatorname{Hom}(}\ \varphi(X)\times\psi(Y)\longrightarrow\operatorname{Hom}(Y,X)\]

% loc: lower argument
Dire que $\beta$ a un adjoint a droite \inserted{(partout défini)} signifie que \deleted{\unsure{pour}} le foncteur contravariant $\psi$ est représentable. Prenant $\varphi(X)=\varnothing$ pour tout $X$, la donnée du morphisme bifonctoriel $\varpi$ devient triviale. Exprimons que $\beta\alpha'$ a un adjoint à droite, i.e. que le foncteur sur $E$

% loc: lower functor
\[X\longmapsto\operatorname{Hom}_{F}(\beta\alpha'(X),a)=\psi(\alpha'(X))\]

% loc: bottom question and reformulation
est représentable. Donc la question devient simplement : Un foncteur $\psi:E_1^{\circ}\to(\mathrm{Ens})$ tel que $\psi\circ\alpha'$ est rep, est-il représentable ? \deleted{\unsure{On peut identifier|Identifions}}\inserted{Identifiant} $\psi$ à un foncteur sur $E^{\circ}$ qui transforme objets de $\Sigma$ en bij, on trouve : Soit \deleted{$\unsure{Z}$} $\xi\in\operatorname{Ob}E$ tel que le \inserted{contr.} foncteur qu'il représente sur $E$, $X\mapsto\operatorname{Hom}_E(X,\xi)$, transforme les $u\in\Sigma$ \deleted{bijections} en bijections. Alors est-il vrai que $\exists\,\xi_1\in E_1$

% archive sheet number: 6
\scanpage{8}

% loc: top, continuation
tel qu'on ait un isom fonctoriel en $X$

% loc: first Hom display
\[\operatorname{Hom}_{E}(X,\xi)\simeq\operatorname{Hom}_{E_1}(\alpha'X,\xi)\quad(\simeq\operatorname{Hom}_{E}(X,\alpha\xi))\]

% loc: question
en d'autres termes a-t-on $\xi\in\operatorname{Ob}\overline{E_1}$ (image essentielle de $\alpha$) ? Or quelle est la relation entre $\Sigma$ et $E_1$ ?

% loc: equivalences
\[
\begin{aligned}
\deleted{\unsure{u\in\Sigma}}\ (u:X\to Y)\in\Sigma
&\Longleftrightarrow\alpha'(u)\inserted{:\alpha'(X)\simeq\alpha'(Y)}\ \text{un isom}
\Longleftrightarrow\underset{\text{pour }\forall\, Z_1\in\operatorname{Ob}E_1}{\operatorname{Hom}_{E_1}(\alpha'(Y),Z_1)\simeq\operatorname{Hom}_{E_1}(\alpha'(X),Z_1)}\\
&\Longleftrightarrow\operatorname{Hom}_{E_1}(Y,\alpha Z_1)\simeq\operatorname{Hom}_{E_1}(X,\alpha Z_1)\quad\text{pour tt } Z_1\in\operatorname{Ob}E_1,
\end{aligned}
\]

% loc: last paragraph
Donc la question est si on récupère bien $\overline{E_1}$ comme sous-catégorie str. pleine de $E$ comme la sous catégorie \inserted{str.} \deleted{\unsure{pleine}} $\Phi(\Sigma)\subset E$ formée des $\xi$ \unsure{tel|tels} que $\operatorname{Hom}(Y,\xi)\xrightarrow{\sim}\operatorname{Hom}(X,\xi)$ pour tt $u:X\to Y$ dans $\Sigma$. Or c'est vrai sauf erreur \ldots.

% archive sheet number: 7
\scanpage{10}

% loc: top, abandoned opening
Soit $A$ une \deleted{$\unsure{\otimes}$}-catégorie \deleted{\unsure{commutative}}\\
\deleted{\unsure{Comme une|une} $\otimes$ catégorie $\unsure{\mathrm{AUC}|\mathrm{ANC}}$ \unsure{est triviale|dont tous} (et définie à isomorphisme \unsure{)}}

% loc: upper paragraph
\deleted{\uline{\unsure{Problème|Définition}}} Soit $\Phi(A)$ la catégorie des familles finies d'objets de $A$, un morphisme de $(I,(L_i)_{i\in I})$ dans $(J,(M_j)_{j\in J})$ étant un couple $(\tau,(\varphi_i)_{i\in I})$ d'une bijection $\tau:I\to J$ et d'isom. $\varphi_i:L_i\xrightarrow{\sim}M_{\tau(i)}$. On introduit sur $\Phi(A)$ une $\otimes$-structure en posant

% loc: tensor formula
\[(I,(L_i)_{i\in I})\otimes(J,(M_j)_{j\in J})=(I\amalg J,(L_i;M_j))\ \ldots\]

% loc: structural constraints
\deleted{\unsure{et}} une structure \deleted{d'associativité}\inserted{d'associativité}, provenant de l'associativité dans l'opération $I\amalg J$, une structure unité par la famille vide \unsure{étant|et avec|étant aussi} l'objet unité, une structure de commutativité par la commutativité de $I\amalg J$. (NB. pour ces définitions, la $\otimes$-structure de $C$ n'a pas encore servi.)

% loc: cancelled functor draft
On \deleted{peut trouver un $\otimes$-foncteur AUC} \deleted{\unsure{pour}}
\[
\deleted{\Pi:\Phi(C)\longrightarrow C\qquad
\substack{\Phi_n(C)=\text{ss-catégorie}\\\text{pleine de }\Phi(C)\text{ formée}}}
\]
\deleted{tel que (la restriction de $\Pi$ au \unsure{le}}\inserted{\deleted{\unsure{notant} $\Pi_n$}}

% loc: canonical functor
a un foncteur canonique
\[\Phi(A)\longrightarrow\big(\underset{\substack{\text{avec isom}\\\text{comme}\\\text{morphismes}}}{\mathrm{Ens\ finis}}\big)\deleted{\unsure{\text{ds Ens}}}\]

% loc: middle, cardinality
(foncteur transportable \ldots), on désigne par $\Phi_n(A)$ l'image inverse \deleted{\unsure{de la}} la sous-catégorie des ensembles de \deleted{\unsure{dim}} cardinal $n$ ($n$ entier $\geq 0$). On a donc

% loc: middle display
\[\Phi(A)=\coprod_n\Phi_n(A),\quad\left\{\begin{aligned}
\Phi_0(A)&=\{e\}\\
\Phi_1(A)&\simeq A
\end{aligned}\right.\]

% loc: 1, statement
1) Soit $C$ une catégorie munie d'un produit tensoriel $\otimes$ AUC. Considérons le foncteur ``restriction à $\Phi_1(A)$''

% loc: restriction display
\[\underline{\operatorname{Hom}}^{\otimes\,\mathrm{AUC}}(\Phi(A),C)\longrightarrow\underline{\operatorname{Hom}}(A,C)\]

% loc: proposition
\uline{Proposition}\quad C'est une équivalence de catégories.

% loc: proof
\uline{Dém} \ldots.

% loc: corollary
\uline{Corollaire}\quad \inserted{Soit $C$ une $\otimes$-catégorie AUC.} Il y a, à isom unique près, un et un seul $\otimes$-foncteur AUC

% loc: bottom
\[\Pi:\Phi(e)\longrightarrow C\qquad\Pi=(\Pi_n)\]

% archive sheet number: 9
\scanpage{11}

% loc: top continuation
munie d'un isom. $\alpha$ rendant commutatif

% loc: top diagram
\diagram{\begin{tikzcd}[ampersand replacement=\&,row sep=small]
\deleted{\unsure{\Phi\,\Pi\,\Phi(e)}}\&\\
\Phi_1(e)\arrow[r,shift left=.9ex,"\Pi_1", ""{name=A,below}]\arrow[r,bend right=16, "{\text{can}^{\unsure{\mathrm{on}|\circ}}}"',""{name=B,above}]\&C\arrow[Rightarrow,from=A,to=B,"\alpha"]\\[-.7em]
\deleted{\unsure{\gamma}}\arrow[u]\arrow[r,dash]\&\deleted{\unsure{e}}
\end{tikzcd}}

% loc: definition
\uline{Définition}\quad Si $(I,(L_i)_{i\in I})$ est une famille finie d'objets de $C$, $\Pi((L_i))$ est appelé le produit $\otimes$ de la famille dans $C$.

% loc: converse
Il y a une réciproque : Si on se donne \deleted{\unsure{une catégorie}} \deleted{\unsure{avec}}\inserted{une cat.} \inserted{avec} une opération $\otimes$, alors la donnée d'une structure $\unsure{\mathrm{AUC}|\mathrm{ANC}}$ sur celle-ci est \unsure{équivalent|équivalente} à la donnée d'un foncteur $\Pi:\Phi(e)\to C$, \deleted{et d'un isom. $\Pi(X\otimes Y)\simeq\Pi(X)\otimes\Pi(Y)$, en associant à tte structure d'associativité} qui sur $\Phi_1(C)$ soit ``l'identité'', et \deleted{\unsure{sur}} qui \deleted{\unsure{sur}} rend commutatif de plus le diagramme

% loc: middle triangle
\diagram{\begin{tikzcd}[ampersand replacement=\&]\Phi_2(C)\arrow[r,"\Pi_2"]\& C\\C\times C\arrow[u]\arrow[ur,"\otimes"']\&\end{tikzcd}}

% loc: ruled separator
\rule{2.3em}{1.4pt}

% loc: NB
\uline{NB}\quad Dans $\Phi(C)$ lui-m\textsuperscript{e}, si on a une famille finie d'objets de $\Phi(C)$, $(\xi_j)_{j\in J}$, $\xi_j=(I_j,(L_{ji})_{i\in I_j})$, alors $\Pi((\xi_j)_{j\in J})$ est

% loc: NB display
\[\left(\coprod_{j\in J}I_j,\ (L_{ij})_{\substack{j\in J\\i\in I_j}}\right)\quad\text{(avec abus de \uline{langage})}\]

% loc: foot paragraph
Revenons à la $\otimes$-\unsure{ct|cat} $C$ AUC. Soit $I$ un ens. \deleted{\unsure{(I)}} à \uline{deux} él\textsuperscript{nts}, $(L_i)_{i\in I}$ une famille d'objets de $C$ indexée par $I$. Il y a deux ordres totaux sur $I$, et si les deux él\textsuperscript{nts} de $I$ sont notés $a,b$, alors les deux façons de représenter le produit, suivant l'ordre total choisi, sont $L_a L_b$ et $L_b L_a$, on passe de l'un à l'autre par l'isom de symétrie. Cela offre \deleted{\unsure{deux}} du nouveau si $L_a=L_b=L$, car alors on trouve un \deleted{\unsure{isom} $s(L)$ \unsure{de L sur lui}} autom. d'ordre 2 de $L\otimes L$ \deleted{\unsure{qui}}

% archive sheet number: 10
\scanpage{12}

% loc: construction
\uline{Contraction}\quad Soit $\Phi^c(C)$ la catégorie formée des \deleted{familles} couples $(\xi,\gamma)$, où $\xi\in\operatorname{Ob}\Phi(\unsure{C|1})$, $\xi=(I,(L_i)_{i\in I})$ (les $L_i\in\operatorname{Ob}C$), et où $\gamma$ est formé

% loc: data, part 1
1) d'un ensemble \inserted{$\gamma'$} (qui peut être vide) de parties de $I$, qui sont \unsure{chacune|de card.} de cardinal 2 et deux à deux disjointes

% loc: data, part 2
2) Pour $\forall J\in\gamma'$, d'un hom de

% loc: pairing
\[
\begin{aligned}\Pi(\xi|J)&\longrightarrow1_C\\
\rotatebox{90}{$=$}&\\
L_{j'}\otimes L_{j''}&\\[-2pt]\text{si }J=\{j',j''\}&
\end{aligned}
\]

% loc: left, beside pairing
% slanted
\margin{$\bigcup\gamma'=\bigcup_{\alpha\in\gamma'}J_\alpha$\\ \unsure{est appelé l'ens} \unsure{des|de}\\ \unsure{ou|des} indices contractés de \deleted{\unsure{$(\Phi^c(C))$}}\inserted{\unsure{$\xi$|$\gamma$}}}

% loc: bracket under pairing
[qui soit un accouplement parfait\\ i.e. une dualité]

% loc: morphisms and restriction
\deleted{\unsure{Les}} Les hom (qui sont des isom) de $\Phi^c(C)$ clairs. On trouve un foncteur évident \inserted{$\Phi^c(C)\to\Phi(C)$} (\unsure{une paire de foncteurs|en plus des foncteurs|des foncteurs} d'inclusion inc : $\operatorname{res}:\Phi^c(C)\longrightarrow\Phi(C)$

% loc: forgetting the pairs
(oubli des indices contractés). $\Phi^c(\unsure{C|1})$ est aussi une $\otimes$-catégorie ass. comm. \deleted{Ceci dit,} unitaire, et res est un $\otimes$ foncteur $\unsure{\mathrm{AUC}|\mathrm{ANC}}$.

% loc: last paragraph
Ceci dit, on trouve un hom canonique

% loc: contraction arrow
\[\Pi\circ\operatorname{inc}\xrightarrow[\text{Contraction}]{c}\Pi\circ\operatorname{res}\]

% loc: bottom
de $\otimes$-foncteurs AUC de $\Phi^c(C)$ dans $C$.

\noindent\hfil\rule{2.7em}{.4pt}\hfil

% archive sheet number: 11
\scanpage{13}

% loc: top continuation
donc il n'y a pas d'isom \uline{canonique} entre $L\otimes L$ et \deleted{\unsure{le}} produit de la famille constante \inserted{sur} $I$ \deleted{\unsure{des|de}} valeur $L$ : il faut, pour choisir un tel isom, choisir un ordre total sur $I$. c'est pire si $I$ est de cardinal $\geq2$, il y a \deleted{\unsure{plus de}} a priori $n!$ isom. possibles entre $\Pi((L_i)_{i\in I})$ et $L^{\otimes 2}$,

% loc: invertible L
Mais si $L$ est inversible, \unsure{la|le} symétrie de $L\otimes L$ \deleted{est} s'identifie à un objet \inserted{$\varepsilon(L)$} du groupe $\operatorname{Aut}(1_C)$, d'ordre 2, et le passage d'un isom. de $\Pi L_i$ avec $L^{\otimes n}$ à un autre, via une permutation $\sigma$ de $[1,n]$, se fait par $\varepsilon(L)^{\operatorname{sign}(\sigma)}$. Il n'y a aucune ambiguité pour le choix dans le seul cas des catégories de Picard [tout objet inv, les $\varepsilon(L)=1$].

% loc: left of inverse discussion
% slanted
\margin{\unsure{kif kif} pour \unsure{pour que l'un des 2|pour un des 2} mais $L\otimes L'\to1$ n'est plus un isom.}

% loc: inverse definition
Inverse d'un élément $L$ inversible : c'est un $L'$ muni d'un isom
\[L\otimes L'\simeq1\]
d'où, en appliquant la symétrie, un \deleted{\unsure{1}} autre isom
\[L'\otimes L\simeq1\]
donc la relation entre $L$ et $L'$ est symétrique, on convient de dire que $(L,L')$ (\deleted{\unsure{sans|non}} ordre fixé) est un couple d'inverses si on se donne un isom
\[\underline{\Pi(L,L')\simeq1}\]

% loc: middle left triangle
\diagram{\begin{tikzcd}[ampersand replacement=\&] L\otimes L'\arrow[d,"s_{L,L'}"']\arrow[dr,"\varphi_{L,L'}"]\&\\L'\otimes L\arrow[r,"\varphi_{L',L}"']\&1\end{tikzcd}}
\uline{\unsure{comm|canon|Comme}}

% loc: middle right
\[L\otimes L'\otimes L\xrightarrow[\sim]{\varphi_{L,L'}\otimes\operatorname{id}_L\simeq\operatorname{id}_L\otimes\psi_{L',L}}L\qquad
\psi_{L',L}:L'\otimes L\simeq1\]
Relation entre $\varphi_{L',L}$ et $\psi_{L',L}$ ? \unsure{Ils|Id.}

% loc: lower formulas
\[
\begin{aligned}
&\varphi_{L,L'}\otimes\operatorname{id}_L\qquad\underbrace{L\otimes L'}\otimes L\longrightarrow L\\
&\operatorname{id}_L\otimes\varphi_{L',L}=\operatorname{id}_L\otimes(\varphi_{L,L'}\circ s_{L',L}):L\otimes L'\otimes L\longrightarrow L
\end{aligned}
\]

% loc: foot left, auxiliary square
$\varphi:M\to1$
\diagram{\begin{tikzcd}[ampersand replacement=\&,column sep=small] L\otimes M\arrow[rr,"s_{L,M}"]\arrow[d,"L\otimes\varphi"']\&\&M\otimes L\arrow[d,"\varphi\otimes L"]\\L\otimes1\arrow[rr,"{\unsure{s|k}_{L,1}}"]\arrow[dr,"\simeq"']\&\&1\otimes L\arrow[dl,"\simeq"]\\\&L\&\end{tikzcd}}

% loc: foot centre, two contractions
\diagram{\begin{tikzcd}[ampersand replacement=\&,column sep=tiny,row sep=small]
\underbrace{L\qquad L'}\arrow[d,"\varphi"]\& L\&\quad L\&\underbrace{L'\qquad L}\arrow[d]\\
1\&\& L\&\underbrace{L^{\unsure{*|a}}\qquad L'^{\unsure{*|a}}}\arrow[d,"\varphi"]\\
\&\&\&1
\end{tikzcd}}

% loc: foot right, bracket and array
\diagram{$\left]\ \begin{array}{cccc}&L&L'&L\\ &L&L^{\unsure{\xi|*}}&L'\\ \unsure{\xi|\varepsilon}&L&L'&L\end{array}\right.$}

% archive sheet number: 12
\scanpage{14}

% loc: top sheet number
\mspage{2}

% loc: top
On a donc

% loc: first square
\diagram{\begin{tikzcd}[ampersand replacement=\&,column sep=large]
\underline{\operatorname{Hom}}_{\unsure{u|{}}}(I;T,T')\arrow[r]\arrow[d]\&\underline{\operatorname{Hom}}(T,T')\arrow[d]\\
{[\underline{\operatorname{Hom}}_0(I;T,T')}\arrow[r]\&{\underline{\operatorname{Hom}}_0(T,T')]}
\end{tikzcd}}

% loc: upper right of square
; avec \unsure{comm|conv}\textsuperscript{on} \unsure{en} \unsure{$\simeq$|---} si $\overline{T},\overline{T'}$ sont de \unsure{prés.} \unsure{$\varprojlim_0$}

% loc: first paragraph
% A single long diagonal stroke traverses the middle of the sheet.
Les théories \deleted{\unsure{lim}}\inserted{fini} $\varprojlim_0$ \inserted{(resp. $\varprojlim$)} forment une 2-catégorie \deleted{\unsure{les|soit}} $\unsure{\mathscr{T}|\tau|\mathcal{T}}$ (resp. $\mathscr{T}_0$), et quand $I$ est fixé, \deleted{\unsure{disons}}\inserted{via les} \deleted{\unsure{ou}} \unsure{$I$|str.}-morphismes de théories, on a \deleted{\unsure{une}} 2-catégories $\mathscr{T}^{(I)}$ ($\mathscr{T}_0^{(I)}$), et \deleted{des} \inserted{un diagramme de} foncteurs de 2-catégories

% loc: second square
\diagram{\begin{tikzcd}[ampersand replacement=\&]
\mathscr{T}\arrow[r]\arrow[d]\&\mathscr{T}^{(I)}\arrow[d,"\deleted{\unsure{\uparrow}}"']\\
\mathscr{T}_0\arrow[r]\&\mathscr{T}_0^{(I)}\end{tikzcd}}

% loc: crossed text at square
\deleted{\unsure{pour qui une fois} $\mathscr{T},\mathscr{T}$}\\
\deleted{Enfin, si on considère les théories qui sont définies par}

% loc: constructions
\uline{Constructions} \inserted{(resp. construction infinie)} pour une \deleted{\unsure{cat.}} théorie $T$ : c'est par définition un hom \deleted{dans} la théorie $i$ \deleted{$\unsure{T_0\ E_0}$} correspondant à un ens de base sans structure, donc $i_C=C$ ; donc c'est un foncteur de \unsure{2|la}-cat. 2-fibrées sur $(\mathrm{Cat}\,\varprojlim\ \text{fin})$ (resp. $(\mathrm{Cat}\,\varprojlim)$)

% loc: construction arrow
\[T\longrightarrow i\qquad\deleted{(\text{resp. }\unsure{T_0}\to i_0)}\]

% loc: category of constructions
Ces constructions \deleted{\unsure{(resp. $T$)}} forment donc une catégorie $\Gamma^T_{\deleted{\unsure{0}}}=\underline{\operatorname{Hom}}(T,i)$ (resp. \inserted{$\Gamma_0^T=$}$\underline{\operatorname{Hom}}_0(T,i)$], qui a des $\varprojlim$ finies (resp. des $\varprojlim$) et qui est munie de

% loc: evaluation
\[\begin{aligned}\Gamma^T\times T&\longrightarrow i\\(\text{resp.}\quad\Gamma_0^T\times T&\longrightarrow i_0)\end{aligned}\]

% loc: last lines
Pour $\forall$ \inserted{$C\in(\mathrm{Cat}\,\varprojlim_0)$ (resp. $\in\mathrm{Cat}\,\varprojlim$) et} $\xi\in T(C)$, le foncteur
\[\begin{aligned}u&\longmapsto u(\xi)\\{\unsure{\Pi|\Gamma}}^{T}&\longrightarrow C\quad(\text{resp. }\Gamma_0^{T}\to \unsure{C|\supset})\end{aligned}\]
commute aux $\varprojlim$ finies \deleted{\unsure{lorsque|précédentes}} (resp. quelc.)

% archive sheet number: 13
\scanpage{16}

% loc: upper question
\uline{Question}\qquad $C$ petite catégorie

$\varphi:C\longrightarrow\mathrm{Ens}$ foncteur exact à droite [$C$ stable par $\varinjlim$ finies] alors $\bar\varphi:(\widehat{C^\circ})^\circ\unsure{\to|{-}}\mathrm{Ens}$ qui prolonge $\varphi$ et commute aux petites $\varprojlim$ (donné aussi par $\bar\varphi(\psi)=\operatorname{Hom}(\psi,\varphi)\quad\psi\in\operatorname{Ob}\underline{\operatorname{Hom}}(C,\mathrm{Ens})$) est-il exact à droite, i.e.

% loc: upper left diagram
\diagram{\begin{tikzcd}[ampersand replacement=\&] {\bar C=(\widehat{C^\circ})^\circ=\underline{\operatorname{Hom}}(C,\mathrm{Ens})^\circ}\arrow[dr]\&\\ C\arrow[u,dash]\arrow[r,"\varphi"]\&\mathrm{Ens}\end{tikzcd}}

% loc: first functor
\[\psi\longmapsto\operatorname{Hom}(\psi,\varphi)\]
transforme-il lim. projectives finies en lim inductives ?

% loc: upper NB
\uline{NB}\quad Par hypoth. sur $\varphi$, c'est le cas pour les $\underset{\longleftrightarrow}{\lim}$ finies qui proviennent de $C$

[\uline{NB} l'inclusion $C\to\bar C$ commute aux $\varinjlim$ \deleted{\unsure{finies}} -- finies ou non --]

% loc: abelian variant
\deleted{\unsure{Ou bien peut supposer que}}\inserted{Supposons} $C$ est \deleted{\unsure{additive}} abélienne, pour $\varphi=$ ou encore $\psi\to I_2^{\operatorname{Hom}(\psi,\varphi)}$ est-il exact à g.
\[\underline{\operatorname{Hom}}(C,\mathrm{Ens})\unsure{\to|{-}}(\mathrm{Ens})\]

% loc: middle, injectivity
$\deleted{\operatorname{Hom}(\psi_1\times\psi_2,\varphi)\unsure{\to|\simeq}}$ Il faudrait que $\varphi$ transforme \unsure{monos|mono} \deleted{\unsure{(fibres)|en}} \inserted{en} épi i.e. $\varphi$ est ``injectif'' :

% loc: extension triangle
\diagram{\begin{tikzcd}[ampersand replacement=\&]\psi\arrow[r,hook]\arrow[dr]\&\psi'\arrow[d,dashed,"\exists"]\\\&\varphi\end{tikzcd}}

% loc: middle consequence
. Cela implique (faire $\psi=\varphi,\ \psi'=\varphi+\rho$) que pour $\forall\,\rho\ \unsure{\in|:}\, C\to\mathrm{Ens}$, on a $\underline{\operatorname{Hom}}(\rho,\varphi)\ne\varnothing$

% loc: left, second triangle
\diagram{\begin{tikzcd}[ampersand replacement=\&]\varphi\arrow[r]\arrow[dr,equal]\&\varphi\amalg\psi\arrow[d,dashed]\\\&\varphi\end{tikzcd}}

% loc: contradiction
\deleted{\unsure{autrement|C'est vrai si}} \uline{Mais c'est tjrs faux} \deleted{sauf si $C$ \unsure{est vide|vide} \illegible{1}} car si on prend pour $\rho$ le foncteur constant de valeur $x$ non vide, on a $\operatorname{Hom}(\rho,\varphi)=\varnothing$ car

% loc: Hom at empty object
\[\operatorname{Hom}_{\mathrm{Ens}}(\rho(\varnothing_C),\varphi(\varnothing_C))=\operatorname{Hom}_{\mathrm{Ens}}(x,\varnothing)=\varnothing\quad !\]

% loc: second question, data
\uline{Question}\qquad Soient $C$ petite catégorie avec $\varprojlim$ finies

$E$ catégorie avec $\varprojlim$ finies et petites $\varinjlim$

$\varphi:C\to E$ un foncteur exact à gauche

$\bar\varphi:\hat C\to E$ son prolongement canonique (commutant aux \inserted{petites} $\varinjlim$)

% loc: second question
Sous quelles conditions d'exactitude sur $E$ peut-on dire que $\bar\varphi$ est aussi exact à g ? [\uline{NB} si $E=(\mathrm{Ens})^\circ$, ce n'est \uline{jamais} le cas !]

% loc: last paragraph
On sait que c'est OK si $E$ est un topos. \uline{Réciproque} ? Supposons que $\varphi$ soit strictement générateur, \deleted{Alors le fait que} alors $E\to\hat C$ pl. \deleted{Il suffit que ce soit un topos.} fidèle et $\bar\varphi:\hat C\to E$ en est un adjoint à g. Donc si $\bar\varphi$ est exact à g, $E$ est un topos.

% loc: bottom right, upside down
% margin note written upside down
\diagram{\begin{tikzcd}[ampersand replacement=\&] \deleted{\unsure{C}/F}\arrow[d,dash]\&\&\\ \deleted{\unsure{C}}\&\deleted{x}\arrow[r]\arrow[d]\&\deleted{F}\\\&\deleted{y}\arrow[ur]\&\end{tikzcd}}

% archive sheet number: 15
\scanpage{17}

% loc: upper right, in reading orientation
% The writing is upside down relative to the archive number. A long pencil stroke crosses the central column.
\margin{$T_p(\Omega^{\unsure{*|\times}})$}

% loc: field diagram
\diagram{\begin{tikzcd}[ampersand replacement=\&,column sep=small]
\mathbb Z_p\arrow[r,hook]\arrow[d,dash]\&W\arrow[r,hook]\arrow[d,dash]\&V\arrow[d,dash]\&\&\\
\mathbb Q_p\arrow[r,dash]\&K\arrow[r,dash]\arrow[rr,bend right=20,dash,"\deleted{\unsure{\pi}}"']\&L\arrow[r,dash]\arrow[r,bend right=22,dash,"\pi"']\&\bar K\arrow[r,dash]\&\Omega
\end{tikzcd}}

% loc: filtered crystals
$M$ \deleted{\unsure{un}} $F$-\unsure{cristalx|cristaux} filtrés sur $(\unsure{k|V},L)$

% loc: upper boxed formulas
\[\deleted{\unsure{M}}\ =\quad\boxed{\boxed{G=\operatorname{Aut}(\bar T_p)}}\qquad\boxed{\boxed{\pi\xrightarrow{\rho}G(\mathbb Q_p)}}\]

% loc: fibre functor
\deleted{$T_p(M)$}\quad $\bar T_p(\underline M)\otimes_{\mathbb Q_p}\Omega$\quad foncteur fibre à val. dans $\Omega$

\deleted{\unsure{avec}} autom \quad $G_\Omega$

% loc: graduation
ce foncteur fibre a une graduation (invariante sous $\pi$ opérant diagonalement) d'où

% loc: boxed homomorphism
\[\boxed{\boxed{\unsure{j|f}_\Omega:\mathbb G_m\longrightarrow G_\Omega}}\]

% loc: conjugation
tel que pour tout $\sigma\in\pi$, on a
\[\boxed{{}^{\sigma}j_\Omega=\operatorname{int}\rho(\sigma)^{-1}\circ j_\Omega}\]

% loc: beside conjugation, left
\margin{$\operatorname{int}\rho(\sigma)\circ{}^{\sigma}j=j$}

% loc: beside conjugation, right
\margin{$\deleted{\unsure{\exists\,M(i)=}}$\\$x\in M_{\Omega}(i)\unsure{\Rightarrow|\to}$\\$\rho(\sigma)\cdot{}^{\sigma}x\in M_{\Omega}(i)$\\i.e. $ {}^{\sigma}x\in\rho(\sigma)^{-1}M_{\Omega}(i)$\\i.e. $ {}^{\sigma}M_{\Omega}(i)=\rho(\sigma)^{-1}M_{\Omega}(i)$}

% loc: expanded formula
\[{}^{\sigma}j_\Omega(\lambda)=\rho(\sigma)^{-1}j_\Omega(\lambda)\rho(\sigma)\]

% loc: lower isomorphism
De plus, on a un $\otimes$-isom \unsure{$K$|$L$}-linéaire \unsure{fonct.|für} $M$ \unsure{variable}
\[
\operatorname{gr}(\bar T_p(M)_\Omega)\simeq\sum_i\operatorname{gr}^i(\mathrm{DR}_{L}(M))\otimes\underbrace{\unsure{\bar T_p|T_p}(T^{\unsure{i|'}}(1))}_{\deleted{\unsure{\text{de Tate}}}\ S_\Omega^{\unsure{i|{}}}}
\]
où

% loc: lower right Tate object
où $T'(1)$ est l'objet de Tate
\[\begin{aligned} S_\Omega&=\bar T(T'(1))\quad\Omega\text{-vectoriel}\\&=\bar T_p(\Omega^{*})^{-1}\end{aligned}\]

% loc: foot left
\[P\simeq P^f\qquad G'\simeq G'^{f}\]

% loc: foot centre
\[\mathbb G_m\xrightarrow{j'_L}G'_L\supset U'_L\]
défini mod conj. par $\in U_L(L)$

% archive sheet number: 16
\scanpage{18}

% loc: header, three columns
\diagram{$\begin{array}{ccc}
C\underset{i}{\subset} & E & F\\
\substack{\text{sous}\\\text{-catégorie}\\\text{pleine}\\{}[\text{strictement}\\\text{génératrice}]}&
\substack{\text{catégorie}\\{}[\text{avec petites}\\\text{limites}\\\text{inductives}]}&
\substack{\text{catégorie avec}\\\text{petites limites}\\\text{inductives}}
\end{array}$}

% loc: pair of functors
\diagram{\begin{tikzcd}[ampersand replacement=\&] {\underline{\operatorname{Hom}}(C,F)}\arrow[r,shift left=.7ex,"\alpha"]\&{\underline{\operatorname{Hom}}(E,F)}\arrow[l,shift left=.7ex,"\beta"]\end{tikzcd}}

% loc: extension and restriction
\[
\left\{\begin{aligned}\alpha(\varphi)(X)&=\varinjlim_{Y\in C/X}\varphi(Y)\\
\beta(\psi)(Y)&=\psi(Y)\qquad\text{i.e. }\beta(\psi)=\psi|C=\psi\circ i
\end{aligned}\right.
\]

% loc: first composite
\[\beta\alpha(\varphi)(\deleted{\unsure{X}}Y)=\alpha(\varphi)(Y)=\varinjlim_{Z\in C/Y}\varphi(Z)=\varphi(Y)\]

% loc: right of first composite
\margin{i.e. $\beta\alpha(\varphi)\simeq\varphi$\\$\beta\alpha\simeq\operatorname{id}_{\underline{\operatorname{Hom}}(C,F)}$}

% loc: second composite
\[\alpha\beta(\psi)(X)=\varinjlim_{Y\in C/X}\beta\psi(Y)=\varinjlim_{Y\in C/X}\deleted{\unsure{\varphi}}\psi(Y)\longrightarrow\psi(X)\]
i.e.

% loc: right of second composite
\margin{$\alpha\beta(\psi)\to\psi$\\i.e. $\alpha\beta\to\operatorname{id}_{\underline{\operatorname{Hom}}(E,F)}$}

% loc: adjunction
\[\operatorname{Hom}(\alpha\varphi,\psi)\simeq\operatorname{Hom}(\varphi,\beta\psi)\]

% loc: before condition
En fait, $\alpha=i_!^F$. Donc $i$ étant pl. fidèle, $i_!$ l'est.

Supposons maintenant que [$E$ soit aussi stable par \inserted{petites} $\varinjlim$] $C$ \inserted{strictement génératrice} et voyons à quelle condition $\alpha(\varphi)$ commute à celles ci. Comme $\beta\alpha(\varphi)=\alpha(\varphi)|C=\varphi$, il est évidemment nécessaire que l'on ait

% loc: starred condition beside vertical stroke
$(*)$ Pour $\forall$ système inductif $X_i$ dans $C$ dont la limite inductive \uline{dans} $E$ est \inserted{un} $Y\in\operatorname{Ob}C$ [à fortiori sa $\varinjlim$ dans $C$ est $Y$] on a
\[\varphi(Y)\xleftarrow{\sim}\varprojlim_i\varphi(X_i)\]

% loc: bottom, subcategories
Soit $\underline{\operatorname{Hom}}'(C,F)$ la sous catégorie pleine de $\underline{\operatorname{Hom}}(C,F)$ formée des $\varphi$ satisfaisant $(*)$ et $\underline{\operatorname{Hom}}'(E,F)$ la sous catégorie pleine de $\underline{\operatorname{Hom}}(E,F)$ formée des $\psi$

% archive sheet number: 17
\scanpage{19}

% loc: foot of scanned image, left in the direction of writing
% The two entries are upside down. The faint reversed writing is show-through from scan 18.
\[\deleted{\varinjlim X_{\unsure{i|'}}}\]

% loc: foot of scanned image, right in the direction of writing
\diagram{\begin{tikzcd}[ampersand replacement=\&]X_i\arrow[r]\& X\\ \&Y\arrow[u]\end{tikzcd}}

% archive sheet number: 18
\scanpage{20}

% loc: top continuation
commutant aux petites $\varinjlim$. Donc $\beta$ induit $\beta':\underline{\operatorname{Hom}}'(E,F)\unsure{-|\to}\underline{\operatorname{Hom}}'(C,F)$. On a pour $\psi$ de $\underline{\operatorname{Hom}}'$

% loc: first calculation
\[\alpha\beta(\psi)(X)=\varinjlim_{Y\in C/X}\psi(Y)\xrightarrow{\sim}\psi(X)\quad\text{car}\quad X=\varinjlim_{Y\in C/X}Y,\ C\]

% loc: faithfulness
étant str. gén. Donc on trouve que $\beta':\underline{\operatorname{Hom}}'(E,F)\to\underline{\operatorname{Hom}}'(C,F)$ est pleinement fidèle. C'est une équivalence i.e. ess. surj. sss $\alpha$ \deleted{\unsure{induit}} envoie bien $\underline{\operatorname{Hom}}'(C,F)$ dans $\underline{\operatorname{Hom}}'(E,F)$. Or soit $(X_i)$ un syst. inductif dans $E$, \deleted{on}\inserted{ayant dans $E$ une limite inductive $Y$} peut alors trouver un système inductif $Y_{ji}$ dans $C$, \inserted{sur l'ens d'indices $I'$ \unsure{est|et}} fibré sur $I$, \unsure{telle|tels|tel} que les \deleted{\unsure{lim}}\inserted{syst. inductifs} partiels (à $i$ fixé) \deleted{\unsure{de $E$}}\inserted{sont} les $C/X_i$, \inserted{de limite $X_i$} et que le système inductif correspondant soit $(Y_{ij})$. On a donc

% loc: iterated colimits
\[\varinjlim_i\deleted{\unsure{\beta}}\alpha(\varphi)(Y_i)=\varinjlim_i\big(\varinjlim_j\varphi(Y_{ij})\big)=\varinjlim_{ij}\varphi(Y_{ij})\underset{\longrightarrow}{\deleted{\unsure{\simeq\lim}}}\varphi(Y)\]

% loc: reduction
et le composé est l'hom dont on veut prouver que c'est un isom. Donc on est ramené à prouver ceci :

% loc: lemma
\uline{Lemme}\quad Si $\varphi\in\underline{\operatorname{Hom}}'(\unsure{C|E},F)$, alors pour $\forall$ système inductif $Y_i$ dans $C$ ayant une limite inductive $Y$ \uline{dans} $E$, on a $\varinjlim\varphi(Y_i)\xrightarrow{\sim}\beta\varphi(Y)$.

% loc: final comparison
Il faut donc comparer \deleted{\unsure{le}} les images par $\varphi$ des 2 systèmes inductifs $(Y_i)$ et $C/Y\to C$. Evidemment on a

% loc: left of cofinality question
\margin{?}

% loc: final question
$I\to C/Y$, cela est-il cofinal ? C'est vrai en tous cas si pour $\forall$ diagramme \uline{fini} dans $C$, la $\varinjlim$ \uline{dans $E$} existe et est dans $C$.

% archive sheet number: 19
